\documentclass[12pt]{article}
\usepackage[utf8]{inputenc}
\usepackage[T1]{fontenc}
\usepackage{hyperref}
\title{Pangu Language Resources for K-12 Curriculum Development}
\author{Compiled by Claude Sonnet 5.5 (claude-sonnet-5-5) for Andr\'as Kornai}
\date{}
\begin{document}
\maketitle

\section{Introduction}
Pangu, usually written Pongu (also called Rin; ISO 639-3 \texttt{png}, Glottocode pong1250), is a Kainji (Atlantic-Congo) language of Niger State, Nigeria, centred on Pangu Gari near Tegina, Rafi LGA. Wikipedia reports about 20,000 speakers in text and 30,000 (2003) in the infobox. There are eight Rin clans with mutually comprehensible dialects. The Endangered Languages catalogue lists the language as vulnerable. Joshua Project says Islam predominates and reports Scripture portions since 2009. The search found no school materials.

\section{Primary Educational Resources}
No primers, textbooks or school programmes were found.

\section{Linguistic and Reference Materials}
\begin{itemize}
\item Wikipedia, ``Pongu language''. \url{https://en.wikipedia.org/wiki/Pongu_language}
\item Endangered Languages Project, Pongu entry and bibliography. \url{https://www.endangeredlanguages.com/lang/2206/bibliography?hl=en}
\end{itemize}

\section{Religious and Community Resources}
\begin{itemize}
\item Joshua Project, Pongu in Nigeria. \url{https://joshuaproject.net/people_groups/14443/NI}
\end{itemize}

\section{Digital and Audio Resources}
\begin{itemize}
\item Global Recordings Network, Pongu (Azhiga) language page. \url{https://glorecnet-website.k.globalrecordings.net/en/language/15737}
\end{itemize}

\section{Curriculum Development Context}
Pongu is not known to be taught in schools. Little written material was found; the bibliography at the Endangered Languages Project is the most promising lead for linguistic documentation. A curriculum would require new orthography and materials development.

\end{document}